\documentclass[12pt]{article}
\usepackage[utf8]{inputenc}
\usepackage[T1]{fontenc}
\usepackage{hyperref}
\title{Fongoro Language Resources for K-12 Curriculum Development}
\author{Compiled by Claude Sonnet 5.5 (claude-sonnet-5-5) for Andr\'as Kornai}
\date{}
\begin{document}
\maketitle

\section{Introduction}
Fongoro (also called Gele; ISO 639-3 fgr) is a nearly extinct language of Chad, formerly also of Sudan. Its classification is uncertain: it is listed as Central Sudanic (Nilo-Saharan), and Glottolog discusses an analysis as a dialectal variant of Gula (kcm). Sources describe it as spoken only by older people, with the community shifting to Fur, Arabic and other regional languages. Search found only encyclopedic and catalogue entries; no teaching materials were found.

\section{Primary Educational Resources}
None found. No primers, textbooks or school programmes for Fongoro appeared in search results.

\section{Linguistic and Reference Materials}
\begin{itemize}
\item Wikipedia, Fongoro language (overview and classification): \url{https://en.wikipedia.org/wiki/Fongoro_language}
\item Glottolog entry for ISO code fgr (endangerment rating, dialect note): \url{https://glottolog.org/resource/languoid/iso/fgr}
\item ASJP database entry (word list data): \url{https://asjp.clld.org/languages/FONGORO}
\item Verbix endangered-languages page: \url{https://docs.verbix.com/EndangeredLanguages/Fongoro}
\end{itemize}

\section{Religious and Community Resources}
Joshua Project has a profile of the Fongoro people in Chad: \url{https://www.joshuaproject.net/index.php/people_groups/11219/CD}. No scripture translations or community organisations were found.

\section{Digital and Audio Resources}
None found.

\section{Curriculum Development Context}
There is essentially no documentation usable for a school curriculum. Any programme would first require fieldwork with the few remaining speakers, documentation and orthography development, and would be revitalisation rather than mother-tongue instruction.

\end{document}